\chapter{課題: 同時観測イベントの探索}
\label{chap:kadai}
\index{どうじかんそく@同時観測|textbf}
\section{入力データと実行環境の準備}

この課題では、次のファイルを配布する。

\begin{itemize}
\item \code{events\_data.zip}
\item \code{slow\_event\_analysis.py}
\item 本テキスト
\end{itemize}

\code{events\_data.zip} を展開すると、\code{events\_data/} というディレクトリが得られる。日ごとに gzip 圧縮されたイベントファイルがあり、年・月ごとのディレクトリに分かれている。

課題のスクリプトやデータは、GitHub リポジトリ
（\url{https://github.com/i-komae/python_textbook}）
から入手できる。手元に配布物がそろっていない場合は、リポジトリの取得と、直下にある \code{slow\_event\_analysis.py} および \code{events\_data.zip} の配置確認が必要である。\code{events\_data.zip} の展開によって得られる \code{events\_data/} が入力ディレクトリとなる。

配布物は意図的に最小限にしてある。したがって、入力、正解判定の対象となるファイル、自分で実装する範囲を着手前に整理する必要がある。

\subsection{配布コードによる基準測定}

最初の基準測定では、配布されたコードを変更せずに実行せよ。

\begin{lstlisting}[numbers=none, xleftmargin=2\zw]
$ python3 slow_event_analysis.py events_data --outdir slow_output
\end{lstlisting}

全期間の処理時間は、データ量、ストレージ、計算機によって変わる。初回の対象は数日分とし、入力形式と出力の確認後に対象期間を広げるのがよい。

最初から完成版を書く必要はない。配布コードの処理時間と出力を変更前に記録しておけば、各修正の効果を基準値との比較によって説明できる。

\subsection{入力ファイルの構造}

解析を始める前に、gzip ファイルを一つ開き、各列の意味を確認せよ。

\begin{lstlisting}[numbers=none, xleftmargin=2\zw]
$ gzip -cd events_data/2024/01/events_2024_01_01.txt.gz | head
\end{lstlisting}

各行の列は次のようになっている。

\begin{lstlisting}[numbers=none, xleftmargin=2\zw]
event_id year month day hhmmss ns detector_id signal
\end{lstlisting}

各列の型と役割を言語化すると、整数として扱う列、時刻を構成する列、識別子としてのみ使う列、物理量として解釈する列を区別できる。この区別が曖昧なままでは、型変換や時刻比較の誤りを招く。

この課題では、2 事象の時刻差を ns 単位まで含めて評価し、\(\lvert \Delta t \rvert \le 250\,\mathrm{ns}\) を満たすものを「同時観測」と定義する。したがって、\code{hhmmss} 列だけを見ればよいわけではなく、\code{ns} 列も含めて時刻を一つの量として組み立てる必要がある。秒境界や日付境界をまたぐ場合にどう振る舞うかも、この定義を満たす形で確認しなければならない。
\index{じこくさ@時刻差|textbf}
\index{じかんまど@時間窓|textbf}

\subsection{小規模入力による検証}

500 日分の全データを処理する前に、数日分を別ディレクトリへコピーし、短時間で再実行できる検証用入力を用意せよ。

\begin{lstlisting}[numbers=none, xleftmargin=2\zw]
$ mkdir -p events_data_small/2024/01
$ cp events_data/2024/01/events_2024_01_01.txt.gz events_data_small/2024/01/
$ cp events_data/2024/01/events_2024_01_02.txt.gz events_data_small/2024/01/
$ cp events_data/2024/01/events_2024_01_03.txt.gz events_data_small/2024/01/
\end{lstlisting}

小さい入力での確認は、速度のためではなく、正しさを素早く確かめるためである。数十秒で終わる条件を作っておけば、修正してすぐ再実行できる。大きなデータを回すのは、その後で十分である。

\begin{lstlisting}[numbers=none, xleftmargin=2\zw]
n_pairs_total = 241
analysis_sec = 18.4
\end{lstlisting}

この段階では、数値そのものよりも、修正のたびに一貫した結果が出るかが大切である。件数が毎回変わる、出力ファイルの形式が揺れる、といった状態では、その先の高速化に進むべきではない。

\section{検証・修正・高速化の要件}

必須項目は次の五点である。

\begin{enumerate}
\item 配布コードの処理時間の測定
\item 同時観測イベント数の検証
\item バグの特定と修正
\item 解析処理の高速化
\item 同時観測イベントの時刻差分布の作図とフィット
\end{enumerate}

この順序にも意味がある。性能改善に着手するのは、元のコードの振る舞いを把握し、正しさを確認した後である。誤ったコードを高速化しても、誤った結果が速く得られるだけである。

\subsection{必須出力}

必須出力は次の三点である。

\begin{itemize}
\item 同時観測イベントの一覧
\item 同時観測イベントを構成する 2 事象の時刻差の分布
\item 同時観測イベントどうしの発生間隔、またはそれに相当する量の分布
\end{itemize}

ここでいう「2 事象の時刻差」は、同時観測と判定した 2 つのイベントの時刻差である。たとえば $\Delta t = t_2 - t_1$ と定義してよい。ただし、符号付きの値と絶対値のどちらを使うか、検出器番号によって向きをそろえるかを定め、その定義を明記せよ。

分布図には、ビン幅、横軸の範囲、単位を明記せよ。図だけを提示しても、量の定義と計数方法が不明では解釈できない。フィットを行う場合は、使用した範囲も記載せよ。

\begin{lstlisting}[numbers=none, xleftmargin=2\zw]
coincidence_pairs.csv
pair_dt_hist.pdf
pair_mean_gap_hist.pdf
summary.txt
\end{lstlisting}

ファイル名の流儀を最初に決めておくと、後で解析を繰り返したときにも混乱しにくい。どのファイルが一覧で、どれが図で、どれが数値要約なのかが名前から分かるようにしておくとよい。

\subsection{正しさを検証する基準値}
\index{ただしさのけんしょう@正しさの検証|textbf}

今回の配布データは 500 日分であり、正しく同時観測イベントを拾えていれば、全期間で得られる同時観測イベント数は

\begin{lstlisting}[numbers=none, xleftmargin=2\zw]
n_pairs_total = 12033
\end{lstlisting}

になる。この値が最終確認の目安である。

ただし、この値に合ったからといって、それだけで実装が正しいとは限らない。偶然件数だけ一致している可能性もある。実装の検証には、可能な範囲で個々のイベント対の内容と、ファイル境界付近の振る舞いも含める必要がある。

\begin{lstlisting}[numbers=none, xleftmargin=2\zw]
n_pairs_total = 12033
analysis_sec = 248.6
plot_sec = 1.8
\end{lstlisting}

このような要約ファイルを出しておくと、コード変更のたびに結果を比較しやすい。件数だけでなく処理時間も同じ場所へ残すと、正しさと速度を一緒に追える。

\section{提出レポートと発展的実装}

提出物は次の二つである。

\begin{itemize}
\item 解析に使った Python コード
\item LaTeX で書いたレポートと、その PDF
\end{itemize}

レポートには、少なくとも次の内容を含めよ。

\begin{itemize}
\item 元のコードの問題点
\item 正しさをどう検証したか
\item バグをどう修正したか
\item どう高速化したか
\item 改善前後の処理時間
\item 作成した図
\item フィット結果とその解釈
\end{itemize}

レポートの構成は、変更点の羅列ではなく、問題の発見、原因の切り分け、修正方針、検証方法、結果、考察の順が望ましい。図表は本文から参照し、その内容と解釈を説明せよ。

\begin{lstlisting}[numbers=none, xleftmargin=2\zw]
1. 元のコードの観察
2. バグの特定
3. 修正内容
4. 高速化方針
5. 結果の比較
6. 図とフィットの考察
\end{lstlisting}

この順序で書けば、読者は「何が問題で、どう直して、何が改善したのか」を追いやすい。レポートは成果物であると同時に、思考過程の記録でもある。

\subsection{発展課題}
\index{べんちまーく@ベンチマーク}

余裕があれば、図の作成やフィットまで含めた全体処理をさらに高速化してよい。さらに、NumPy や pandas を使うだけでなく、自分でアルゴリズムを設計し直したり、C や C++ で計算コアを実装したりしてもよい。

この発展課題では、特定の解法を要求しない。pandas を深く使う方向もあれば、ストリーミング処理を突き詰める方向もあり得るし、重い部分だけを別言語へ切り出す方向もあり得る。重要なのは、自分の方針を説明できることである。
